\documentclass{article}
\usepackage{mathtools} 
\usepackage{fontspec}
\usepackage[UTF8]{ctex}
\usepackage{amsthm}
\usepackage{mdframed}
\usepackage{xcolor}
\usepackage{amssymb}
\usepackage{amsmath}


% 定义新的带灰色背景的说明环境 zremark
\newmdtheoremenv[
  backgroundcolor=gray!10,
  % 边框与背景一致，边框线会消失
  linecolor=gray!10
]{zremark}{说明}

% 通用矩阵命令: \flexmatrix{矩阵名}{元素符号}{行数}{列数}
\newcommand{\flexmatrix}[4]{
  \[
  #1 = \begin{pmatrix}
    #2_{11}     & #2_{12}     & \cdots & #2_{1#4}   \\
    #2_{21}     & #2_{22}     & \cdots & #2_{2#4}   \\
    \vdots      & \vdots      & \ddots & \vdots     \\
    #2_{#31}    & #2_{#32}    & \cdots & #2_{#3#4}
  \end{pmatrix}
  \]
}

% 简化版命令（默认矩阵名为A，元素符号为a）: \quickmatrix{行数}{列数}
\newcommand{\quickmatrix}[2]{\flexmatrix{A}{a}{#1}{#2}}

\begin{document}
\title{3.2 注释}
\author{张志聪}
\maketitle

\begin{zremark}
  对定理3.4直观版本：

  定理 3.4$^\prime$： 若$\lim\limits_{x \to x_0} f(x) = A$，那么对任意实数
  $B < A < C$，存在$\delta > 0$，当$0 < |x - x_0| < \delta$时，有
  \begin{align*}
    B < f(x) < C
  \end{align*}
\end{zremark}

\textbf{证明：}

$\lim\limits_{x \to x_0} f(x) = A$，由定义
对$\epsilon = min\{A - B, C - A\}$，$\exists \delta > 0$，
当$0 < |x - x_0| < \delta$时，有
\begin{align*}
  |f(x) - A| < \epsilon \\
  B \leq A - \epsilon < f(x) < A + \epsilon \leq C
\end{align*}

\begin{zremark}[复合函数极限定理]
  设函数 $f$ 与 $g$ 满足：
  \[
    \lim_{x \to u_0} f(x) = a, \quad \lim_{t \to t_0} g(t) = u_0,
  \]
  并且当 $t$ 在 $t_0$ 的某去心邻域内时，$g(t)$ 的值都属于 $f$ 的定义域。
  则复合函数 $f(g(t))$ 在 $t_0$ 处的极限存在，且有
  \[
    \lim_{t \to t_0} f(g(t)) = a.
  \]
\end{zremark}

以上，也是“变量替换法（换元法）”在极限应用的理论依据。
\end{document}